\section{Pre-existing Notetaking Software}
\label{research:ExistingNotetakingSoftware}
\textbf{Source}: Useful opinions from industry experts\cite{wvfrmpodcastOurFavoriteProductivity2024}
\subsection*{OneNote}

OneNote has been my note storge system for many years thanks to its robust integration with other Microsoft tools and reliable (if not the most feature rich) collaboration tools . While it has many drawing features and some organisation tools (notebooks, sections and subsections) these often become frustrating to use especially when trying to link larger projects together (such as this EPQ). This has been and will continue to be (at least until I develop my software) my primary note taking resource, therefore I have a page above entirely dedicated to noting down frustrations I face using it in my everyday life.

The main pro of OneNote (its design being built of Microsoft office) is also its fundamental flaw as it is restricted to the same file system used since the 1980s and the same corporate aesthetics that aren't always suitable for personal learning resources.


\begin{tabularx}{\textwidth}{>{\raggedright\arraybackslash}X>{\raggedright\arraybackslash}X}
    \rowcolor{gray!15} \textbf{Pros}                                         & \textbf{Cons}                                                                                                             \\
    Robust drawing tools                                                     & Lacking organisation features                                                                                             \\
    Excellent typing tools                                                   & Some tools feel clunky and poorly implemented                                                                             \\
    Image / pdf compatibility (including very helpful "print to pdf" feature & Notes often appear dull and monotonous, one part of which is due to the ugly default colour scheme across office products \\
    Collaboration (especially teacher - student)                             &                                                                                                                           \\
    Cross compatibility                                                      &                                                                                                                           \\
\end{tabularx}

\subsection*{Notability}
For many reasons notability has been my go-to for producing aesthetically pleasing, handwritten notes on my iPad. It combines amazing drawing tools with a slick interface. The onboard pastel colour palette for example is very easy on the eyes and helps direct attention to different resources.

It does not however, support reliable cross compatibility and collaboration as files often fail to sync between my devices and is restricted to only my Apple devices (no windows support which Is frustrating).

\begin{tabularx}{\textwidth}{>{\raggedright\arraybackslash}X>{\raggedright\arraybackslash}X}
    \rowcolor{gray!15} \textbf{Pros} & \textbf{Cons}                                                             \\ Amazing drawing tools   &   Poor collaboration features \\
    Good typing tools                & Poor cross platform compatibility                                         \\
                                     & Lacking organisation features                                             \\
                                     & Strange issues occasionally such as lost data or unmodifiable text fields \\
\end{tabularx}

\subsection*{Obsidian}
Obsidian is a hugely powerful note storage tool which can be customised to fit nearly every use case. It allows for note tagging, subsections and so much more. Whilst these tools are incredible powerful they are also incredibly time consuming to setup and manage as your note library grows. For this reason it has seen little mass market appeal outside of further education users and those who don’t mind learning the various scripting features necessary to make the most of the tool.

As one of those users I love obsidian for its flexibility and whilst I know I cannot match its flexibility (as this also brings complexity which I want to avoid), I aim to make it much simpler to use the key tools and feature it has and in which I have identified as important through my study of the human learning process.

Unfortunately obsidian doesn’t have many tools towards handwritten notes and diagrams which has meant I cant use it for subjects such as maths, where handwritten notes are often far easier, and more powerful than typed alternatives (such as LaTeX).

\begin{tabularx}{\textwidth}{>{\raggedright\arraybackslash}X>{\raggedright\arraybackslash}X}
    \rowcolor{gray!15} \textbf{Pros}           & \textbf{Cons}                                                             \\ Extremely robust organisation tools   &   Difficult to setup and huge learning curve \\
    Very customisable (even more so with code) & Poor cross platform compatibility                                         \\
                                               & Very limited handwritten note support and information recorded in written form cant be integrated with most of the powerful storing and searching tools obsidian had                                           \\
                                              
\end{tabularx}

\subsection*{TickTick}
Tick Tick isn't a notetaking app in the same way as the others but more so a diary / reminder system but the features within it are very powerful. Tick Tick allows you to create reminders using computationally process natural language, for example typing "Friday morning take out bins" would create a new reminder automatically set for Friday morning (the exact reminder time taking into account your past choices) and tagged / searchable by any key term. This idea of natural language interaction really appeals to me and in envisage my tools scanning all text and assigning tags and other ideas to aid searchability.

Additionally the tool has advanced calendar / diary tools which I want to replicate within my artifact to help users track changes to files over time - after all time was one of the searching methods I found Iny my searching research. 

\begin{tabularx}{\textwidth}{>{\raggedright\arraybackslash}X>{\raggedright\arraybackslash}X}
    \rowcolor{gray!15} \textbf{Pros} & \textbf{Cons}                                                             \\ Very natural and fluid interaction  &    \\
  Powerful date / time integration               &                         \\
                                    
\end{tabularx}

\subsection*{Apple Notes}
Apple notes is a very basic notetaking tool which allows you to store notes in folders, whilst its feature set is very generic (and limited) I find the concept of quick notes (notes created without any organisation but that can be later expanded into main notes very useful. This helps me note quick ideas without needing to worry about organising or losing data.

\begin{tabularx}{\textwidth}{>{\raggedright\arraybackslash}X>{\raggedright\arraybackslash}X}
    \rowcolor{gray!15} \textbf{Pros} & \textbf{Cons}                                                             \\ Notes can always be found  &  Whilst can become main notes, the main note system itself is not particularly advanced / noteworthy compared to OneNote, obsidian etc\\
    Able to draw, write and attach files / images             &                                  \\
            Can set alarms to remind myself to revisit notes                         &                        \\
                                    
\end{tabularx}

\subsection*{Apple FreeForm}
Freeform is a really new and interesting take on creating collaborative notes. Documents are simply a huge (near infinite) canvas in which text, drawings, images and shapes can be added and seen by all. Sections of ideas can be managed by views in which locations around the whiteboard can be recorded and jumped to by participants.

Unfortunately this is largely the extent of the organisation features within freeform, you can create multiple boards but this feels 'hacky' as it is not clear if you should use one huge board for everything and jump around or split ideas around.

\begin{tabularx}{\textwidth}{>{\raggedright\arraybackslash}X>{\raggedright\arraybackslash}X}
    \rowcolor{gray!15} \textbf{Pros} & \textbf{Cons}                                                             \\ Infinite canvas for ideas  &  Only available on apple ecosystem devices limiting collaboration \\
    Reliable collaboration tools              & Finding specific notes within the canvas can be very frustrating as there is no organisation feature beyond views           \\
                                     & Collaboration options only exist for the board as a whole, specific sections cant be shared which can be frustrating if one idea (which should have one board) has to be split to preserve privacy                                           \\
                                  
\end{tabularx}


\subsection*{Zotero}
I discovered Zotero during a homework tasks and while it isn't exactly designed for general note taking, it does have very powerful filtering and tagging features which I can learn from. I downloaded the software and gave it a try by filling it in with my current source tracker.

The auto adding features of books and research papers is very flexible and I would like to integrate a similar solution into my software.

\begin{tabularx}{\textwidth}{>{\raggedright\arraybackslash}X>{\raggedright\arraybackslash}X}
    \rowcolor{gray!15} \textbf{Pros} & \textbf{Cons}                                                             \\ Flexible tagging system that allows unlimited tags per resource  &  Application feels quite corporate and designed for academic / research use which might not appeal to younger students \\
    Automatic resource adding based on public data (like ISBN)             & Whilst basic tags exist they can't be chained as parent child relationships which makes the creation of many specialised tags which Is frustrating                            \\
                                  
\end{tabularx}
